%% ============================================================
%%  entorno.tex  —  Sección 1.2
%% ============================================================

\subsection{Entorno experimental}

\subsubsection{Hardware}

Los experimentos se ejecutaron en el equipo local del equipo de desarrollo
bajo el sistema operativo \textbf{Windows} (versión detectada durante la
ejecución de los benchmarks compilados con MinGW/MSYS2).

\begin{table*}[t]
  \centering
  \caption{Especificaciones de hardware del entorno experimental.}
  \label{tab:hardware}
  \begin{tabular}{ll}
    \toprule
    \textbf{Componente} & \textbf{Especificación} \\
    \midrule
    Procesador     & \texttt{Intel(R) Core(TM) i5-1035G1} \\
    Frecuencia     & \texttt{CPU @ 1.00GHz base} \\
    Núcleos / Hilos & \texttt{Cores: 4, LogicalProcessors: 8, MaxClockSpeed: 1190} \\
    Memoria RAM    & \texttt{16\,GB (16553776Kb)} \\
    Almacenamiento & \texttt{Model: KINGSTON SNV2S500G, Size: 500105249280} \\
    \bottomrule
  \end{tabular}
\end{table*}

Los benchmarks se ejecutan de forma \textbf{monohilo} (proceso único sin
paralelismo explícito). La medición de tiempo utiliza
\lstinline{std::chrono::high_resolution_clock} de la biblioteca estándar de
C++, lo que introduce una resolución típica de microsegundos en plataformas
Windows modernas.

\subsubsection{Software}

\begin{table*}[t]
  \centering
  \caption{Software utilizado en el proyecto.}
  \label{tab:software}
  \begin{tabular}{lll}
    \toprule
    \textbf{Componente} & \textbf{Herramienta} & \textbf{Versión} \\
    \midrule
    Compilador C++ & \texttt{g++} (MinGW-w64 / MSYS2) &
      \texttt{14.2.0} \\
    Estándar C++   & C++17 & — \\
    Flags de compilación & \texttt{-O2 -std=c++17 -Wall} & — \\
    Python         & CPython & 3.13\footnotemark \\
    Biblioteca \texttt{pandas}   & pandas  & 2.3.2 \\
    Biblioteca \texttt{matplotlib} & matplotlib & 3.10.6 \\
    Biblioteca \texttt{seaborn}  & seaborn & 0.13.2 \\
    Control de versiones & Git & \texttt{git version 2.43.0.windows.1} \\
    \bottomrule
  \end{tabular}
\end{table*}

\footnotetext{Versión inferida de la ruta de instalación
\texttt{C:\textbackslash{}Users\textbackslash{}...\textbackslash{}Python313}
observada durante la instalación de dependencias.}

Las versiones de \texttt{pandas}, \texttt{matplotlib} y \texttt{seaborn}
fueron registradas automáticamente durante la ejecución de
\texttt{pip install} en el entorno del equipo de desarrollo.

Las optimizaciones del compilador (\texttt{-O2}) se aplicaron en todos los
experimentos para reproducir condiciones de compilación representativas del
software de producción, evitando artificios de código no optimizado.
La bandera \texttt{-Wall} se utilizó para capturar advertencias del
compilador sin suprimir ninguna.
